\documentclass{article}

\newcounter{listing}

\usepackage{fontspec}
\setmainfont{DejaVu Serif}
\usepackage[russian,english]{babel}
\usepackage{amsthm}

% ==============================================================================
% LOAD REFSUITE PACKAGE WITH LUA CONFIGURATION
% ==============================================================================
\usepackage{refsuite}
\LoadLuaReferencesConfiguration{lua-flow-conf.lua}

\newtheorem{theorem}{Теорема}

\begin{document}

\title{RefSuite Lua Configuration Flow}
\author{John Doe}
\date{\today}
\maketitle

\tableofcontents

\section{Introduction}
\label{sec:introduction}

This example demonstrates the Lua configuration flow for \texttt{refsuite}. The
package reads a Lua table and uses it to configure short references,
bibliographies, and indexes before the document starts.

The Lua table and bibliography files live beside this source file to keep the
example explicit and easy to inspect.  Build this file from the \texttt{examples}
directory; auxiliary files go to \texttt{.build/}, and the resulting PDF stays
beside this source file.

\subsection{Package Configuration Snippet}

The package is loaded normally, then configured from a Lua file:

\begin{quote}
\begin{verbatim}
\usepackage{refsuite}
\LoadLuaReferencesConfiguration{lua-flow-conf.lua}
\end{verbatim}
\end{quote}

The Lua file can describe multiple bibliography groups and named index commands:

The bibliography path is \texttt{.} because the \texttt{.bib} files are stored
next to this example source.  This keeps the data files visible while still
sending temporary build output to \texttt{.build/}.

\begin{quote}
\begin{verbatim}
return {
  hyperref = {
    enable = true,
    shortref = {
      { csname = 'picref', reftext = 'Fig.' },
      { csname = 'tabref', reftext = 'Table' },
      { csname = 'secref', reftext = 'Section' },
      { csname = 'lstref', reftext = 'Listing' },
    },
  },
  zrefclever = {
    enable = true,
    setup = { lang = 'russian' },
    azcref = true,
  },
  bibtex = {
    enable = true,
    loadtime_opts = { backend = 'biber', style = 'numeric' },
    path = '.',
    bibliographies = {
      { name = 'main', title = 'References', files = 'lua-main.bib' },
      { name = 'online', title = 'Online Resources', files = 'lua-online.bib' },
    },
  },
  indexes = {
    enable = true,
    list = {
      { name = 'notation', title = 'List of Notations', intoc = true },
      { name = nil, title = 'Subject Index', intoc = true },
    },
  },
}
\end{verbatim}
\end{quote}
\refstepcounter{listing}\label{lst:configuration}

\section{Main Content}

\subsection{Configured Short References}

The Lua configuration creates commands such as \picref{fig:sample},
\tabref{tab:sample}, \secref{sec:conclusion}, and \lstref{lst:configuration}.
These commands are generated from the Lua table before the document body begins,
so they can be used like ordinary LaTeX commands.

\begin{figure}[htbp]
  \centering
  \fbox{\parbox{0.5\textwidth}{\centering Example figure}}
  \caption{Example figure with caption}
  \label{fig:sample}
\end{figure}

\begin{table}[htbp]
  \centering
  \begin{tabular}{|c|c|}
    \hline
    Column 1 & Column 2 \\
    \hline
    Data 1 & Data 2 \\
    Data 3 & Data 4 \\
    \hline
  \end{tabular}
  \caption{Example table}
  \label{tab:sample}
\end{table}

\subsection{Bibliographies}

Lua configuration can create named bibliography printing commands. This example
cites books \cite{knuth1984texbook,lamport1994latex} and online resources
\cite{torvalds2003git,ctan2024latex}.

The two bibliography groups are printed separately at the end of the document
using commands generated from their Lua names.

\subsection{Indexes}

The Lua table configures a named notation index and the default subject index.
Both are populated here and printed after the bibliographies.

\index[notation]{$\alpha$} Alpha notation
\index[notation]{$\beta$} Beta notation
\index{LaTeX} Typesetting system
\index{bibliography} List of references
\index{Lua configuration} Lua configuration

\section{Automatic zref-clever References}

The Lua flow can also load \texttt{zref-clever} and define \verb|\azcref| for
automatic Russian case selection when LuaLaTeX is used.

The examples below place different Russian prepositions immediately before
\verb|\azcref|.  The Lua helper inspects the previous words and selects the case
argument passed to \verb|\zcref|.

\begin{otherlanguage}{russian}
\begin{theorem}\zlabel{thm:lua-auto-case}
Тестовое утверждение для проверки автоматического выбора падежа.
\end{theorem}

Без предлога: \azcref{thm:lua-auto-case}.

по \azcref{thm:lua-auto-case}.

к \azcref{thm:lua-auto-case}.

в силу \azcref{thm:lua-auto-case}.

из-за \azcref{thm:lua-auto-case}.
\end{otherlanguage}

\section{Conclusion}
\label{sec:conclusion}

The Lua flow is useful when configuration should live in a Lua data structure or
when named bibliography commands should be generated from that structure.

\printbibliographymain
\printbibliographyonline

\printindexnotation
\printindex

\end{document}
